\section{Virtual DOM}
\label{section:Virtual DOM}

Jest to technologia wykorzystywana przez rózne biblioteki i frameworki (w tym przez React) do efektywnego zarządzania aktualizacjami interfesju w aplikacjach webowych. DOM to Document Object Model, czyli obiektowy model dokumentu. Jego struktura to drzewo, która reprezentuje elementy HTML na aplikacji. Każda zmiana w rzeczywistym DOM jest kosztowna względem wydajności, dlatego Virtual DOM jest niezwykle przydatne w tym przypadku, ponieważ co 200ms ładuje się 800 nowych próbek na wykresy.
\subsection{Działanie Virtual DOM}
Virtual DOM jest uproszczoną wersją drzewa DOM, która jest tworzona przez React i utrzymywana w pamięci. Umożliwia to minimalizowanie liczby bezpośrednich operacji na rzeczywistym DOM. Po każdej zmianie stanu aplikacji Virtual DOM jest porównywany z jego poprzednią wersją, a następnie React dokonuje niezbędnych aktualizacji rzeczywistego DOM, co przekłada się na lepszą optymalizację.

\subsection{Zalety Virtual DOM}
Virtual DOM posiada kluczowe korzyści, które w efektywny sposób przyczyniają się di wydajności i łatwości zarządzania aplikacjami:

\begin{enumerate}
    \item \textbf{Wyższa wydajność}: dzięki aktualizacjom dokonywanym najpierw Virtual DOM, a następnie w rzeczywistym DOM, liczba zasobożernych operacji jest w znacznym stopniu redukowana.
    \item \textbf{Optymalizacja zmian}: proces porównywania zmian w stanie aplikacji pozwala Reactowi dokonywania tylko niezbędnych zmian w rzeczywistym DOM, co znacząco wpływa na wydajność aplikacji. Gdyby za każdym razem stare próbki miały się przeładowywać, to nie byłoby możliwe z dużą częstotliwością przetwarzać dane na wykresach.
    \item \textbf{Prostsze zarządzanie stanem}: dzięki Virtual DOM React może efektywnie zarządzać stanem komponentów, ponieważ aktualizację są dokonywane w myśl sposobu deklaratywywnego, a nie imperatywnego.
    \item \textbf{Modularność}: Virtual DOM umożliwia tworzenie aplikacji z małych i izolowanych komponentów, które w przypadku zmiany stanu będą się aktualizować bez wpływu na inne części aplikacji.
\end{enumerate}
% \subsection{react-plotlyjs}
% React-plotly.js to biblioteka umożliwiająca integrację biblioteki Plotly.js z Reactem. Służy ona do wizualizacji danych z aplikacjami stworzonymi w React. Umożliwia ona tworzenie zaawansowanych i interaktywnych wykresów, które można w bardzo prosty sposób zarządzać i dynamicznie zmieniać wyświetlane dane na wykresie za pomocą React.
% \subsection{Jak działa react-plotlyjs?}
% React Plotly.js oferuje komponent Reactowy, który opiera się na Plotly.js. Wykorzystuje właściwości props do definiowania danych, wyglądu i funkcji interaktywnych wykresów. Dzięki tym narzędziom można w prosty sposób sposób zmieniać dane jak i wygląd wykresów. W aplikacji jest możliwość zmiany jednostek, wspomniana wcześniej biblioteka w prosty sposób umożliwia modyfikacji opisów osi jak i zamianę danych na wykresie.
% \subsection{Zalety react-plotlyjs}
% \begin{enumerate}
%     \item \textbf{Interaktywność:} Zapewnia interaktywne funkcje, takie jak powiększanie, przesuwanie, tooltipy i zapis wykresu.
%     \item \textbf{Integracja z React:} Dzięki komponentowi plotly.js można w prosty sposób zarządzać dynamicznie stanem wykresów. W tym przypadku cały czas wykresy są aktualizowane z aktualną pozycją rotorów jak i uchybem.
%     \item \textbf{Wydajność:} W aplikacji jest wgrywanych, na pojedyńczy wykres, co sekundę aż 1 000 nowych próbek, co w znaczny sposób obciąża aplikację. Jednym z ze sposobów polepszenia wydajności jest zastosowanie WebGL, poprzez typ wykresu scatterg, co znacząco polepsza wydajność przy renderowaniu wykresów.
%     \item \textbf{Eksport danych:} Umożliwia eskport wykresów w formacie PNG.
% \end{enumerate}
% Standardowe wykresy rozproszone w Plotly.js (\textit{scatter}) wykorzystują SVG do renderowania punktów i linii, co jest efektywne dla małych i średnich zbiorów danych. Na wykresach jest ładowana duża ilość próbek, która może być liczona w milionach przez co w tym przypadku renderowanie za pomocą SVG jest niewydajne. W takich przypadkach wykorzystuje się scattergl, który wykorzystuje WebGL.
% Jednym z minusów scattergl jest mniejsza ostrość, ponieważ wykres nie jest renderowany wektorowo.
